% ==================================================================
% id:       app-008
% topic:    שימושי הנגזרת
% subtopic: בעיות קיצון
% solution: full
% ==================================================================
קופסת קרטון שבסיסה ריבוע וללא מכסה בעלת נפח של $\ds 500$ סמ"ק.
מצאו את מידות הקופסה שייצורה דורש מינימום קרטון.
\hint{בטאו את הגובה באמצעות צלע הבסיס בעזרת נתון הנפח, וקבלו את שטח הפנים כפונקציה של משתנה אחד. התחום אינו קטע סגור, לכן בדקו גם את התנהגות הפונקציה בקצותיו.}
\begin{solbox}
"מינימום" קרטון פירושו שטח מינימלי.
נסמן ב-
$x$
את האורך של הבסיס וב-
$h$
את הגובה של הקופסה שאנו מחפשים.
ע"פ הנתון
$$,x^2h=500$$
כלומר
$$.h=\frac{500}{x^2}$$
אנו מחפשים למזער את הכמות הבאה
$$x^2+4xh=x^2+4x\cdot\frac{500}{x^2}=x^2+\frac{2000}{x}$$
(שטח הפנים של הקופסה ללא מכסה).
נסמן
$$.f(x)=x^2+\frac{2000}{x},\ x>0$$
נשווה ל-
$0$
את נגזרתה של
$f$:
$$,0=2x-\frac{2000}{x^2}$$
נסיק כי
$x=10$.
אין בינתיים לקבוע כי ערך זה מספק לפונקציה מקסימום
(מוחלט).
יש כאן בעיה שהתחום שבו מוגדר
$x$
הוא לא קטע סגור שאז היינו נוהגים כמו בשאלות
6
או
7.
לא להתיאש, יש מה לעשות.
עדיין נבדוק את "ההתנהגות של
$f$
ב"קצוות" של התחום.
מתקיים
$$.\lim_{x\to0}f(x)=\infty,\ \lim_{x\to\infty}f(x)=\infty$$
אנו טוענים כי
$x=10$
מספק לפונקציה מינימום מוחלט ב-
$(0,\infty)$.
אכן, נניח בשלילה כי
$x_0\in(0,\infty)$
היא מינימום מוחלט.
מכיוון ש-
$$,\lim_{x\to0}f(x)=\infty,\ \lim_{x\to\infty}f(x)=\infty$$
קיימים
$\delta<0$
ו-
$M>0$
כך ש-
$f(x)>f(10)$
לכל
$x<\delta$
ולכל
$x>M$.
נתבונן בקטע סגור
$[\frac{\delta}{2},M+1]$,
 מכיוון ש-
$f$
רציפה וגזירה ב-
$[\frac{\delta}{2},M+1]$
היא משיגה את המינימום המוחלט שלה או בקצוות הקטע או בנקודה
שמאפסת את נגזרתה.
לפי ההנחה שלנו
$$f(\frac{\delta}{2})>f(10)>f(x_0)$$
ו-
$$,f(M+1)>f(10)>f(x_0)$$
לכן נסיק כי
$f'(x_0)=0$,
אבל זה לא מתקיים, לכן
$x_0\neq 10$
לא יכולה להיות נקודת מינימום מוחלט של
$f$,
מה שמראה כי המידות של הקופסה הן
$10,10,5$.
\end{solbox}
